%==============================================================================
%  Competitive Programming Notebook — Guidebook
%  Autor: Juan  ·  Universidad del Bosque
%  Formato: doble columna (izquierda / derecha), estilo team notebook ICPC
%  Compilar con: pdflatex guidebook.tex   (dos veces, para el índice)
%==============================================================================
\documentclass[10pt,a4paper,twocolumn]{article}

% ---- Codificación e idioma -------------------------------------------------
\usepackage[T1]{fontenc}
\usepackage[utf8]{inputenc}
\usepackage{lmodern}
\usepackage{microtype}
\usepackage{float}
\usepackage{enumitem}
\usepackage{amsmath}   % \text en modo matematico
% Nombres en español (sin depender de babel-spanish, que no está instalado aquí)
\renewcommand{\contentsname}{Índice}
\renewcommand{\partname}{Parte}

% ---- Diseño de página ------------------------------------------------------
\usepackage[a4paper,margin=1.6cm,columnsep=0.8cm]{geometry}
\setlength{\columnseprule}{0.4pt}             % línea fina que divide izquierda/derecha
\usepackage{xcolor}
\definecolor{acento}{RGB}{20,45,85}           % azul marino formal
\definecolor{acentoClaro}{RGB}{70,110,170}
\definecolor{grisCodigo}{RGB}{245,245,247}
\definecolor{comentario}{RGB}{95,120,95}
\definecolor{palabraClave}{RGB}{30,70,150}
\definecolor{cadena}{RGB}{150,80,40}
\def\columnseprulecolor{\color{gray!40}}

% ---- Encabezados y pie -----------------------------------------------------
\usepackage{fancyhdr}
\usepackage{titlesec}
\pagestyle{fancy}
\fancyhf{}
\fancyhead[L]{\small\itshape Competitive Programming Notebook}
\fancyhead[R]{\small\nouppercase{\itshape\rightmark}}
\fancyfoot[C]{\small\thepage}
\renewcommand{\headrulewidth}{0.4pt}

\titleformat{\section}
  {\normalfont\large\bfseries\color{acento}}{\thesection}{0.6em}{}
\titleformat{\subsection}
  {\normalfont\normalsize\bfseries\color{acentoClaro}}{\thesubsection}{0.5em}{}
\titleformat{\part}
  {\normalfont\Large\bfseries\color{acento}\filcenter}{\partname~\thepart}{0.5em}{}

% ---- Resaltado de código ---------------------------------------------------
\usepackage{listings}
\lstset{
  backgroundcolor=\color{grisCodigo},
  basicstyle=\ttfamily\scriptsize,
  keywordstyle=\color{palabraClave}\bfseries,
  commentstyle=\color{comentario}\itshape,
  stringstyle=\color{cadena},
  numbers=left, numberstyle=\tiny\color{gray}, numbersep=5pt,
  showstringspaces=false, breaklines=true, tabsize=2,
  frame=single, rulecolor=\color{gray!40},
  columns=fullflexible, keepspaces=true,
  extendedchars=true,
  literate=
    {á}{{\'a}}1 {é}{{\'e}}1 {í}{{\'i}}1 {ó}{{\'o}}1 {ú}{{\'u}}1
    {Á}{{\'A}}1 {É}{{\'E}}1 {Í}{{\'I}}1 {Ó}{{\'O}}1 {Ú}{{\'U}}1
    {ñ}{{\~n}}1 {Ñ}{{\~N}}1 {¿}{{?`}}1 {¡}{{!`}}1 {°}{{\textdegree}}1
}
\lstdefinestyle{cpp}{language=C++}
\lstdefinestyle{java}{language=Java}
\lstdefinestyle{py}{language=Python}

% ---- Enlaces (índice navegable) --------------------------------------------
\usepackage[colorlinks=true,linkcolor=acento,urlcolor=acentoClaro]{hyperref}

\setcounter{tocdepth}{2}
\setcounter{secnumdepth}{2}

% Marcador de secciones aún sin contenido
\newcommand{\pendiente}{\textit{\color{gray}Contenido en desarrollo.}\par\medskip}

%==============================================================================
\begin{document}

% ===== PORTADA =============================================================
\onecolumn
\thispagestyle{empty}
\begin{titlepage}
  \centering
  \vspace*{\fill}

  {\large\scshape Universidad del Bosque}\\[4pt]
  {\normalsize Facultad de Ingeniería}\\[2pt]
  {\normalsize Ingeniería de Software \& Ingeniería en Robótica}\\[24pt]

  \color{acento}\rule{\linewidth}{1.2pt}\\[8pt]
  {\Huge\bfseries Competitive Programming}\\[10pt]
  {\Huge\bfseries Notebook}\\[10pt]
  {\large\itshape Algoritmos, estructuras de datos y plantillas}\\[8pt]
  {\large C\texttt{++} \quad}\\[6pt]
  \color{acento}\rule{\linewidth}{1.2pt}\\[28pt]

  \color{black}
  {\large\bfseries Juan Carlos Morales}\\[4pt]
  {\normalsize Doble titulación: Sistemas \& Robótica}\\[36pt]

  {\normalsize Bogotá, Colombia \quad$\cdot$\quad 2026}

  \vspace*{\fill}
\end{titlepage}

% ===== ÍNDICE ==============================================================
\thispagestyle{empty}
\tableofcontents
\clearpage

% ===== CUERPO (doble columna) ==============================================
\twocolumn

%---------------------------------------------------------------------------
\part{Fundamentos Básicos}
%---------------------------------------------------------------------------

\section{Entrada / Salida Rápida (Fast I/O)}
Usar lectura rápida para evitar TLE, hay 3 niveles de lectura rápida, estandarizar Tier 2 y usar Tier 3 (lector por
bytes) solo cuando haya del orden de $10^6$--$10^7$ lecturas.

\begin{table}[h]
\centering
\caption{Casos donde se recomienda Fast I/O}
\begin{tabular}{|p{0.36\columnwidth}|p{0.50\columnwidth}|}
\hline
\textbf{Tipo de problema} & \textbf{Ejemplo} \\
\hline
Arrays gigantes & Leer/ordenar/sumar $10^6$--$10^7$ enteros \\
\hline
Muchos test cases (T grande) & $T \leq 10^5$, cada uno trivial pero se acumula \\
\hline
Grafos grandes & $M \sim 10^6$--$2*10^6$ aristas para BFS/DFS/Dijkstra/MST \\
\hline
Geometria con muchas consultas & $L,S$ grandes + muchos test cases seguidos hasta EOF \\
\hline
Muchos tokens de texto & Strings/lineas que suman millones de caracteres \\
\hline
Judges clasicos "trampa de I/O" & SPOJ INTEST, ARRAYSUB, etc. \\
\hline
\end{tabular}
\end{table}

Regla practica: si N/M es del orden de $10^4-10^5$, cin con
\texttt{sync\_with\_stdio(false);} cin.tie(nullptr); alcanza de sobra. El FastReader
es la "artilleria pesada" para cuando eso deja de ser suficiente.

\textbf{Tier 1 — Básico:} \texttt{cin} / \texttt{cout}.
\begin{lstlisting}[style=cpp]
int entero; double decimal; string palabra, linea;
cin >> entero >> decimal >> palabra;  // >> lee UN token
// ! cin >> deja el '\n' pendiente;
// ! el 1er getline sale VACÍO -> hay que ignorarlo
cin.ignore();
getline(cin, linea);                  // línea completa
\end{lstlisting}
\textbf{Tier 2 — Rápido:} desactiva la sync con C (una vez en
\texttt{main}). Punto dulce para casi todo.
\begin{lstlisting}[style=cpp]
ios_base::sync_with_stdio(false);
cin.tie(nullptr);            // ! no mezclar con scanf/printf
cout << entero << "\n";      // "\n", nunca endl (endl frena)
\end{lstlisting}
\textbf{Tier 3 — Más rápido:} lector por bytes con
\texttt{fread}. Solo con $10^6$–$10^7$ lecturas.
\begin{lstlisting}[style=cpp]
struct FastReader {
    static const int TAM = 1 << 16;   // 64 KB
    char bufer[TAM];
    int puntero = 0, leidos = 0;

    int leer() {
        if (puntero == leidos) {
            leidos = (int) fread(bufer, 1, TAM, stdin);
            puntero = 0;
        }
        return leidos == 0 ? -1 : bufer[puntero++]; // -1=EOF
    }
    int nextInt() {
        int resultado = 0, b = leer();
        while (b <= ' ') b = leer();      // saltar blancos
        bool negativo = (b == '-');
        if (negativo) b = leer();
        while (b >= '0' && b <= '9') {
            resultado = resultado * 10 + (b - '0'); b = leer();
        }
        return negativo ? -resultado : resultado;
    }
    long long nextLong() {
        long long resultado = 0; int b = leer();
        while (b <= ' ') b = leer();
        bool negativo = (b == '-');
        if (negativo) b = leer();
        while (b >= '0' && b <= '9') {
            resultado = resultado * 10 + (b - '0'); b = leer();
        }
        return negativo ? -resultado : resultado;
    }
    double nextDouble() {
        double resultado = 0, divisor = 1; int b = leer();
        while (b <= ' ') b = leer();
        bool negativo = (b == '-');
        if (negativo) b = leer();
        while (b >= '0' && b <= '9') {
            resultado = resultado * 10 + (b - '0'); b = leer();
        }
        if (b == '.') {
            b = leer();
            while (b >= '0' && b <= '9') {
                resultado += (b - '0') / (divisor *= 10); b = leer();
            }
        }
        return negativo ? -resultado : resultado;
    }
    string next() {                       // un token
        string s; int b = leer();
        while (b <= ' ') b = leer();
        while (b > ' ') { s += (char) b; b = leer(); }
        return s;
    }
    string nextLine() {                   // línea completa
        string s; int b = leer();
        while (b != '\n' && b != -1) {
            if (b != '\r') s += (char) b; b = leer();
        }
        return s;
    }
};
\end{lstlisting}

\section{Formateo de Salida}


La función \texttt{printf} utiliza una cadena de formato para indicar cómo
debe interpretarse y mostrarse cada argumento. Cada especificador comienza
con \texttt{\%} y corresponde al argumento que aparece después de la cadena.

\begin{lstlisting}[style=cpp]
// %d  -> entero (int)
printf("%d\n", 42);
// Entrada: 42
// Salida: 42

// %lld -> entero largo (long long)
long long x = 1234567890123LL;
printf("%lld\n", x);
// Entrada: 1234567890123
// Salida: 1234567890123

// %5d -> ancho minimo de 5 caracteres, alineado a la derecha
printf("%5d\n", 42);
// Entrada: 42
// Salida: "   42"

// %-5d -> ancho minimo de 5 caracteres, alineado a la izquierda
printf("%-5d|\n", 42);
// Entrada: 42
// Salida: "42   |"

// %05d -> ancho 5, completando con ceros
printf("%05d\n", 42);
// Entrada: 42
// Salida: "00042"

// %+d -> muestra siempre el signo
printf("%+d\n", 42);
// Entrada: 42
// Salida: "+42"

// %x -> hexadecimal usando letras minusculas
printf("%x\n", 255);
// Entrada: 255
// Salida: ff

// %X -> hexadecimal usando letras mayusculas
printf("%X\n", 255);
// Entrada: 255
// Salida: FF

// %.2f -> numero real con exactamente 2 decimales
printf("%.2f\n", 3.14159);
// Entrada: 3.14159
// Salida: 3.14

// %.0f -> numero real sin decimales (redondea)
printf("%.0f\n", 3.7);
// Entrada: 3.7
// Salida: 4

// %e -> notacion cientifica
printf("%e\n", 31415.9);
// Entrada: 31415.9
// Salida: 3.141590e+04

// %c -> un solo caracter
printf("%c\n", 'A');
// Entrada: 'A'
// Salida: A

// %s -> cadena de caracteres
printf("%s\n", "hola");
// Entrada: "hola"
// Salida: hola

// Se pueden combinar varios especificadores
printf("%d %lld %.2f\n", 42, 10000000000LL, 3.14159);
// Entrada: 42, 10000000000, 3.14159
// Salida: 42 10000000000 3.14

// %% -> imprime un '%' literal
printf("100%%\n");
// Entrada: no recibe ningun argumento para %
// Salida: 100%

// ! NO uses %n en C/C++ (puede ser peligroso)
// Para scanf, un double se lee con %lf
\end{lstlisting}

\subsubsection*{Estructura de un especificador}

Un formato puede combinar varias partes:

\begin{lstlisting}[style=cpp]
//   % [flags] [width] [.precision] [specifier]
//   %    0       5         .2          f

printf("%05d\n", 42);
//  %  -> comienza el formato
//  0  -> rellena con ceros
//  5  -> ancho minimo de 5 caracteres
//  d  -> interpreta el argumento como int
//
//  Salida: 00042
\end{lstlisting}

Los modificadores más utilizados en programación competitiva son:

\begin{center}
\begin{tabular}{|c|c|c|}
\hline
\textbf{Formato} & \textbf{Tipo} & \textbf{Ejemplo de salida} \\
\hline
\texttt{\%d} & \texttt{int} & \texttt{42} \\
\hline
\texttt{\%lld} & \texttt{long long} & \texttt{1234567890123} \\
\hline
\texttt{\%c} & \texttt{char} & \texttt{A} \\
\hline
\texttt{\%s} & \texttt{string/char*} & \texttt{hola} \\
\hline
\texttt{\%f} & \texttt{double} & \texttt{3.140000} \\
\hline
\texttt{\%.2f} & \texttt{double} & \texttt{3.14} \\
\hline
\texttt{\%x} & \texttt{int} & \texttt{ff} \\
\hline
\texttt{\%X} & \texttt{int} & \texttt{FF} \\
\hline
\end{tabular}
\end{center}

\textbf{Regla práctica:} en programación competitiva, los formatos más
importantes son \texttt{\%d} para \texttt{int}, \texttt{\%lld} para
\texttt{long long}, \texttt{\%c} para caracteres, \texttt{\%s} para cadenas
y \texttt{\%.2f} (o una precisión diferente) para números reales.


\section{Condicionales y Ciclos}

\begin{lstlisting}[style=cpp]
// if / else: ejecuta un bloque segun una condicion
if (x > 0) cout << "pos";       // x > 0 -> "pos"
else if (x < 0) cout << "neg";  // x < 0 -> "neg"
else cout << "cero";             // x == 0 -> "cero"

// Operador ternario: condicion ? valor_si_true : valor_si_false
string s = (x >= 0) ? "no neg" : "neg";
// Si x >= 0, s = "no neg"; si no, s = "neg"

// switch: compara una expresion con diferentes valores
// Solo funciona directamente con tipos integrales como int o char,
// NO con string.
switch (x) {
    case 1: /*...*/ break;       // x == 1
    case 2: case 3: /*...*/ break; // x == 2 o x == 3
    default: /*...*/ ;           // ningun case coincide
}
// ! Sin break, se continua ejecutando el siguiente case (fall-through)

// C++17: se puede declarar una variable dentro del if
if (int r = x % 2; r == 0) {}
// r solo existe dentro del if y sus bloques asociados

// for clasico: usa un indice y se repite mientras i < n
for (int i = 0; i < n; i++) {}

// range-based for: recorre directamente los elementos
// No proporciona el indice.
for (int v : datos) {}

// & obtiene una referencia al elemento original.
// Por tanto, las modificaciones afectan al vector.
for (int &v : datos) v *= 2;

// while: ejecuta mientras la condicion sea verdadera
while (c < n) c++;

// do-while: primero ejecuta el bloque y DESPUES verifica la condicion.
// Por eso se ejecuta AL MENOS una vez.
do {
    c--;
} while (c > 0);

// break -> termina inmediatamente el ciclo o switch
// continue -> salta a la siguiente iteracion del ciclo

// ! C++ NO tiene break con etiqueta.
// Para salir de ciclos anidados se puede usar una bandera,
// encapsular la logica en una funcion y usar return, etc.
\end{lstlisting}




\section{Conversiones, Parseos y Casteos}

\begin{lstlisting}[style=cpp]
// string -> numero (parseo)
// Convierte una cadena de texto en el tipo numerico indicado.
int e = stoi("42");
// Entrada: "42"  ->  Salida: 42

long long l = stoll("10000000000");
// Entrada: "10000000000"  ->  Salida: 10000000000

double d = stod("3.14");
// Entrada: "3.14"  ->  Salida: 3.14

// Tambien se puede indicar la base numerica.
int bin = stoi("1010", nullptr, 2);
// "1010" en base 2 -> 10

int hex = stoi("ff", nullptr, 16);
// "ff" en base 16 -> 255


// numero -> string
// Convierte cualquier numero compatible en texto.
string s = to_string(42);
// Entrada: 42  ->  Salida: "42"


// Casteo: fuerza la conversion de un tipo a otro.
// Al pasar de decimal a entero, se TRUNCA hacia 0.
int t = (int) 3.99;
// 3.99 -> 3

int r = static_cast<int>(3.99);
// 3.99 -> 3
// static_cast<int> es la forma recomendada en C++


// ! int / int = division ENTERA
double mal = 7 / 2;
// 7 y 2 son int -> 3
// Luego 3 se convierte a double -> 3.0

double bien = 7 / 2.0;
// 2.0 es double -> division decimal
// 7 / 2.0 -> 3.5


// caracteres
// Los char tienen un valor numerico asociado (ASCII).
int cod = (int) 'A';
// 'A' -> 65

char ch = (char) 66;
// 66 -> 'B'

// Truco muy usado para convertir un digito char a int.
int dig = '7' - '0';
// '7' tiene valor 55 y '0' tiene valor 48
// 55 - 48 -> 7
\end{lstlisting}




%---------------------------------------------------------------------------
\part{Estructuras de Datos}
%---------------------------------------------------------------------------

\section{Basic Structures}

\begin{lstlisting}[style=cpp]
int a[5] = {1,2,3,4,5}; // crear con valores: O(n)

int b[5];               // crear de tamano 5
                        // posiciones: b[0] ... b[4]
                        // ! Valores no inicializados

b[0] = 10;              // INSERT: O(1)
b[1] = 20;              // INSERT: O(1)

int x = b[0];           // READ: O(1)

b[0] = 30;              // UPDATE: O(1)

for (int i = 0; i < 5; i++)
    cout << b[i] << " ";// READ: recorrer -> O(n)

// DELETE: no existe metodo para eliminar
// Se puede sobrescribir o manejar logicamente
b[0] = 0;               // DELETE logico: O(1)

// ! El tamano de un array estatico no cambia.
// Para mas posiciones se necesita otro arreglo.
\end{lstlisting}

\subsection{Matrices (Arreglos 2D)}

Arreglo bidimensional (filas $\times$ columnas). \textbf{Complejidad:}
acceso \texttt{[i][j]} $O(1)$; crear $n\times m$ en $O(nm)$.

\begin{lstlisting}[style=cpp]
vector<vector<int>> m(3, vector<int>(4,0)); // CREATE: O(n*m)
// matriz 3x4, todas las posiciones inicializadas en 0

m[1][2] = 7;              // INSERT: O(1)

int x = m[1][2];           // READ: O(1)

m[1][2] = 10;              // UPDATE: O(1)

for (int i = 0; i < 3; i++)
    for (int j = 0; j < 4; j++)
        cout << m[i][j] << " ";
// READ: recorrer toda la matriz -> O(n*m)

// DELETE: no existe metodo para eliminar una posicion.
// Se puede sobrescribir el valor o manejarlo logicamente.
m[1][2] = 0;               // DELETE logico: O(1)

// ! El tamano de la matriz puede cambiar con vector,
// pero un arreglo 2D estatico no puede cambiar de tamano.
\end{lstlisting}



\subsection{Vectores / Listas dinámicas}

Arreglo que \textbf{crece} solo (ArrayList / \texttt{vector} /
\texttt{list}). \textbf{Complejidad:} acceso por índice $O(1)$; agregar
al final $O(1)$ amortizado; insertar/borrar en medio $O(n)$; buscar
$O(n)$.

\begin{lstlisting}[style=cpp]
vector<int> v;              // CREATE: lista vacia

v.push_back(3);             // INSERT: final -> O(1) amortizado
v.push_back(7);             // INSERT: final -> O(1) amortizado

int x = v[0];               // READ: acceso -> O(1)
int n = v.size();            // READ: tamano -> O(1)

v[0] = 10;                  // UPDATE: modificar -> O(1)

v.insert(v.begin() + 1, 5);  // INSERT: medio -> O(n)

v.erase(v.begin() + 1);      // DELETE: medio -> O(n)

// Buscar un elemento recorriendo el vector
for (int x : v)
    if (x == 10) { /* ... */ }
// SEARCH: recorrer -> O(n)

// DELETE: eliminar el ultimo elemento
v.pop_back();               // DELETE: final -> O(1)

// ! vector mantiene acceso por indice O(1).
// ! Al crecer puede realocar su memoria internamente.
\end{lstlisting}



\subsection{Listas enlazadas (LinkedList)}

Nodos enlazados. \textbf{Complejidad:} insertar/eliminar en los
\textbf{extremos} $O(1)$; acceso por índice $O(n)$; buscar $O(n)$. Útil
como cola o deque.

\begin{lstlisting}[style=cpp]
list<int> l;                // CREATE: lista vacia

l.push_front(1);            // INSERT: inicio -> O(1)
l.push_back(9);             // INSERT: final -> O(1)

int x = l.front();          // READ: primer elemento -> O(1)
int y = l.back();           // READ: ultimo elemento -> O(1)

l.front() = 5;              // UPDATE: primer elemento -> O(1)
l.back() = 10;              // UPDATE: ultimo elemento -> O(1)

l.pop_front();              // DELETE: inicio -> O(1)
l.pop_back();               // DELETE: final -> O(1)

// ! No existe acceso directo por indice.
// Avanzar hasta una posicion requiere recorrer nodos.
auto it = l.begin();
advance(it, 2);             // acceso al tercer elemento -> O(n)

// Buscar un elemento requiere recorrer la lista.
for (int x : l)
    if (x == 10) { /* ... */ }
// SEARCH: recorrer -> O(n)

// ! list no permite l[i].
// Para acceso por indice, vector suele ser mejor.
\end{lstlisting}

\subsection{Pilas (Stack)}

Estructura \textbf{LIFO} (último en entrar, primero en salir).
\textbf{Complejidad:} insertar, eliminar y consultar el tope $O(1)$. Para DFS
iterativo, paréntesis balanceados y deshacer.

\begin{lstlisting}[style=cpp]
stack<int> pila;            // CREATE: pila vacia

pila.push(3);               // INSERT: O(1)
pila.push(7);               // INSERT: O(1)

int tope = pila.top();      // READ: consultar tope -> O(1)

pila.top() = 10;            // UPDATE: modificar tope -> O(1)

pila.pop();                 // DELETE: eliminar tope -> O(1)

// ! Solo se puede acceder directamente al elemento superior.
// ! No existe acceso por indice.
\end{lstlisting}


\subsection{Colas (Queue)}

Estructura \textbf{FIFO} (primero en entrar, primero en salir).
\textbf{Complejidad:} encolar y desencolar $O(1)$. Base del BFS.

\begin{lstlisting}[style=cpp]
queue<int> cola;             // CREATE: cola vacia

cola.push(3);                // INSERT: final -> O(1)
cola.push(7);                // INSERT: final -> O(1)

int frente = cola.front();   // READ: primero -> O(1)
int atras = cola.back();     // READ: ultimo -> O(1)

cola.front() = 10;           // UPDATE: modificar primero -> O(1)

cola.pop();                  // DELETE: primero -> O(1)

// ! push agrega al final y pop elimina del frente.
// ! No existe acceso directo por indice.
\end{lstlisting}


\subsection{Colas de prioridad (Priority Queue / Heap)}

Devuelve siempre el elemento de mayor (o menor) prioridad.
\textbf{Complejidad:} insertar y extraer $O(\log n)$; consultar el tope
$O(1)$. Para Dijkstra, Prim, k-ésimo mayor.

\begin{lstlisting}[style=cpp]
priority_queue<int> pq;      // CREATE: MAX-heap por defecto

pq.push(3);                  // INSERT: O(log n)
pq.push(7);                  // INSERT: O(log n)

int mx = pq.top();           // READ: mayor -> O(1)

pq.pop();                    // DELETE: mayor -> O(log n)

// min-heap: el menor queda en el tope
priority_queue<int, vector<int>, greater<int>> pqMin;

pqMin.push(3);               // INSERT: O(log n)
int mn = pqMin.top();        // READ: menor -> O(1)
pqMin.pop();                 // DELETE: O(log n)

// ! No existe acceso directo a elementos intermedios.
// ! priority_queue no tiene una operacion UPDATE directa.
\end{lstlisting}

\subsection{Vector circular (Circular Buffer)}

\textbf{Idea:} un arreglo de tamaño fijo (\texttt{cap}) que se comporta
como una \textbf{cola doble}. En lugar de mover elementos cuando se
inserta o elimina, se reutilizan las posiciones del arreglo de forma
circular.

Se mantienen tres variables principales:

\begin{itemize}
    \item \texttt{cap}: capacidad total del arreglo.
    \item \texttt{ini}: posición física donde comienza el primer elemento.
    \item \texttt{tam}: cantidad actual de elementos almacenados.
\end{itemize}

La posición lógica $i$ no coincide necesariamente con la posición física
del arreglo. Se obtiene mediante:

\[
\texttt{posicionFisica(i) = (ini + i) \% cap}
\]

Por ejemplo, con $\texttt{cap}=5$ y $\texttt{ini}=3$:

\[
\begin{array}{c|ccccc}
\text{Índice lógico} & 0 & 1 & 2 & 3 & 4 \\ \hline
\text{Índice físico} & 3 & 4 & 0 & 1 & 2
\end{array}
\]

Cuando se llega al final del arreglo, el módulo hace que la siguiente
posición vuelva al comienzo. Por eso \textbf{no es necesario mover los
elementos}.

\textbf{¿Por qué usarlo?} En un arreglo normal, insertar al frente
requiere desplazar todos los elementos y cuesta $O(n)$. En un vector
circular solo se mueve \texttt{ini}, por lo que cuesta $O(1)$.

\textbf{Casos de uso:}
\begin{itemize}
    \item Colas de tamaño fijo.
    \item Buffers de audio, video o red.
    \item Ventanas deslizantes.
    \item Sistemas donde llegan y salen datos continuamente.
    \item Simulaciones de procesos en una cola.
    \item Mantener únicamente los últimos $k$ elementos.
\end{itemize}

Si el buffer está lleno y se ejecuta \texttt{pushBack}, el nuevo elemento
\textbf{reemplaza al más antiguo}. Esto permite mantener automáticamente
una ventana de tamaño fijo.

\textbf{Complejidad:} insertar, eliminar, consultar extremos, acceder
por posición y comprobar el tamaño son $O(1)$.


\begin{lstlisting}[style=cpp]
struct VectorCircular {
    vector<int> buf;
    int cap, ini = 0, tam = 0;

    VectorCircular(int c) : buf(c), cap(c) {}

    // CREATE: crea un buffer con capacidad fija.
    // Los elementos inicialmente no estan ocupados.
    
    bool empty() {
        return tam == 0;
    }                           // O(1)

    bool full() {
        return tam == cap;
    }                           // O(1)

    int size() {
        return tam;
    }                           // O(1)

    int capacity() {
        return cap;
    }                           // O(1)

    // Convierte una posicion logica en fisica.
    int pos(int i) {
        return (ini + i) % cap;
    }                           // O(1)

    // INSERT: agrega al final.
    void pushBack(int x) {
        buf[(ini + tam) % cap] = x;

        if (tam < cap)
            tam++;
        else
            ini = (ini + 1) % cap;
    }                           // O(1)

    // INSERT: agrega al inicio.
    void pushFront(int x) {
        ini = (ini - 1 + cap) % cap;
        buf[ini] = x;

        if (tam < cap)
            tam++;
    }                           // O(1)

    // READ: consulta el primer elemento.
    int front() {
        return buf[ini];
    }                           // O(1)

    // READ: consulta el ultimo elemento.
    int back() {
        return buf[(ini + tam - 1) % cap];
    }                           // O(1)

    // READ: acceso por posicion logica.
    int at(int i) {
        return buf[(ini + i) % cap];
    }                           // O(1)

    // UPDATE: modifica un elemento.
    void set(int i, int x) {
        buf[(ini + i) % cap] = x;
    }                           // O(1)

    // DELETE: elimina el primer elemento.
    int popFront() {
        int v = buf[ini];
        ini = (ini + 1) % cap;
        tam--;
        return v;
    }                           // O(1)

    // DELETE: elimina el ultimo elemento.
    int popBack() {
        tam--;
        return buf[(ini + tam) % cap];
    }                           // O(1)

    // DELETE: elimina todos los elementos.
    void clear() {
        ini = 0;
        tam = 0;
    }                           // O(1)

    // SEARCH: busca un elemento.
    // A diferencia de las operaciones anteriores,
    // hay que recorrer los elementos.
    bool contains(int x) {
        for (int i = 0; i < tam; i++)
            if (at(i) == x)
                return true;

        return false;
    }                           // O(n)

    // READ: recorrer los elementos en orden logico.
    void print() {
        for (int i = 0; i < tam; i++)
            cout << at(i) << " ";
    }                           // O(n)
};
\end{lstlisting}

\textbf{Ejemplo de funcionamiento:}

\begin{lstlisting}[style=cpp]
VectorCircular v(5);

v.pushBack(10);
v.pushBack(20);
v.pushBack(30);
// [10, 20, 30, _, _]
// ini = 0, tam = 3

v.pushFront(5);
// [5, 10, 20, 30, _]
// ini cambia de 0 a 4

cout << v.front();
// 5

cout << v.back();
// 30

v.popFront();
// elimina 5
// [_, 10, 20, 30, _]

v.pushBack(40);
v.pushBack(50);
// [_, 10, 20, 30, 40, 50]
// El almacenamiento puede haber dado la vuelta.

v.pushBack(60);
// El buffer estaba lleno.
// Se agrega 60 y se elimina el mas antiguo.
// Resultado logico: [20, 30, 40, 50, 60]
\end{lstlisting}

\textbf{Ventaja principal:} el arreglo físico puede verse ``desordenado'',
pero los elementos siempre tienen un orden lógico determinado por
\texttt{ini}. Por ejemplo:

\begin{lstlisting}[style=cpp]
// Memoria fisica:
// [40, 50, 20, 30,  _]
//       ^
//      ini = 2

// Orden logico:
// 20 -> 30 -> 40 -> 50
\end{lstlisting}

No importa dónde estén físicamente los elementos: \texttt{at(0)} siempre
representa el primero, \texttt{at(1)} el segundo, y así sucesivamente.

\textbf{Diferencia con un vector normal:}

\begin{center}
\small
\begin{tabular}{|l|c|c|}
\hline
\textbf{Operación} & \texttt{vector} & \textbf{Circular} \\
\hline
Agregar al final & $O(1)$* & $O(1)$ \\
\hline
Eliminar al final & $O(1)$ & $O(1)$ \\
\hline
Agregar al frente & $O(n)$ & $O(1)$ \\
\hline
Eliminar al frente & $O(n)$ & $O(1)$ \\
\hline
Acceso por índice & $O(1)$ & $O(1)$ \\
\hline
Buscar & $O(n)$ & $O(n)$ \\
\hline
\end{tabular}
\end{center}

\noindent
*Amortizado.

\textbf{Regla práctica:} si necesitas una estructura de tamaño fijo
donde los elementos entran y salen continuamente por los extremos,
un vector circular evita los desplazamientos de un arreglo normal y
mantiene las operaciones principales en $O(1)$.

\section{Estructuras Hash/Tree}


\subsection{Conjuntos hash (HashSet)}

Colección de elementos \textbf{sin duplicados} y sin orden garantizado.
En C\texttt{++} se implementa mediante \texttt{unordered\_set}. Utiliza
una tabla hash, por lo que las operaciones son $O(1)$ en promedio.
Es útil cuando solo interesa saber si un elemento existe, sin necesidad
de mantener los elementos ordenados.

\textbf{Complejidad:} insertar, buscar y borrar $O(1)$ promedio
($O(n)$ peor caso).

\begin{lstlisting}[style=cpp]
unordered_set<int> s;          // CREATE: conjunto vacio

s.insert(3);                   // INSERT: O(1) prom.
s.insert(7);                   // INSERT: O(1) prom.
s.insert(3);                   // no se agrega: ya existe

bool hay = s.count(3);         // SEARCH: O(1) prom.
// hay = 1 si existe, 0 si no existe

auto it = s.find(7);            // SEARCH: O(1) prom.
// devuelve iterador al elemento o s.end()

s.erase(3);                    // DELETE: O(1) prom.

int n = s.size();              // READ: cantidad -> O(1)
bool vacio = s.empty();        // READ: ¿esta vacio? -> O(1)

// Recorrer todos los elementos
for (int x : s)
    cout << x << " ";
// O(n), pero el orden NO esta garantizado.

// ! No existe acceso por indice: s[0] no es valido.
// ! Los elementos repetidos se ignoran.
// ! El orden puede cambiar entre ejecuciones.
\end{lstlisting}

\textbf{Cuándo usarlo:}
\begin{itemize}
    \item Saber rápidamente si un elemento ya apareció.
    \item Eliminar duplicados.
    \item Mantener un conjunto de elementos visitados.
    \item Detectar elementos repetidos en $O(n)$ promedio.
    \item BFS/DFS: guardar vértices visitados cuando sea conveniente.
\end{itemize}


\subsection{Mapas / Diccionarios hash}

Estructura que almacena pares \textbf{clave $\rightarrow$ valor}.
En C\texttt{++} se implementa mediante \texttt{unordered\_map}.
Cada clave identifica un valor y \textbf{no puede repetirse}. Si se
asigna un nuevo valor a una clave existente, se reemplaza el anterior.

\textbf{Complejidad:} insertar, acceder, buscar y borrar por clave
$O(1)$ promedio ($O(n)$ peor caso).

\begin{lstlisting}[style=cpp]
unordered_map<string,int> m;   // CREATE: mapa vacio

m["a"] = 1;                    // INSERT: O(1) prom.
m["b"] = 2;                    // INSERT: O(1) prom.

int v = m["a"];                // READ: acceder -> O(1) prom.
// v = 1

m["a"] = 10;                   // UPDATE: O(1) prom.
// la clave "a" ahora tiene valor 10

m.erase("a");                  // DELETE: O(1) prom.

bool existe = m.count("b");    // SEARCH: O(1) prom.
// existe = 1

auto it = m.find("b");         // SEARCH: O(1) prom.
if (it != m.end())
    cout << it->second;        // valor asociado a "b"

// Recorrer todas las parejas
for (auto [clave, valor] : m)
    cout << clave << " " << valor;
// O(n), sin orden garantizado.
\end{lstlisting}

\textbf{Importante:} usar \texttt{m[clave]} para consultar una clave
que no existe puede \textbf{crear esa clave} con su valor por defecto.
Para comprobar existencia sin insertar, es preferible usar
\texttt{count()} o \texttt{find()}.

\begin{lstlisting}[style=cpp]
unordered_map<string,int> m;

int x = m["a"];
// Si "a" no existia, se crea:
// "a" -> 0
// ! operator[] puede INSERTAR una clave al consultar.

// count(): comprueba si existe SIN insertar
bool existe = m.count("b");
// Devuelve 0 si no existe, 1 si existe.

// find(): busca la clave SIN insertar
auto it = m.find("b");

if (it != m.end())
    cout << it->second;       // valor asociado a "b"
else
    cout << "No existe";

// ! count() solo indica si existe.
// ! find() permite acceder a la pareja clave -> valor.
\end{lstlisting}


\textbf{Uso clásico: contar frecuencias.}

\begin{lstlisting}[style=cpp]
unordered_map<int,int> freq;

for (int x : datos)
    freq[x]++;                 // frecuencia de cada elemento

cout << freq[5];
// cantidad de veces que aparece 5
\end{lstlisting}

\textbf{Cuándo usarlo:}
\begin{itemize}
    \item Contar frecuencias.
    \item Asociar identificadores con información.
    \item Buscar rápidamente un valor a partir de una clave.
    \item Problemas de conteo y frecuencia.
    \item Memorizar resultados de estados o funciones.
\end{itemize}


\subsection{Conjuntos ordenados (TreeSet)}

Conjunto de elementos \textbf{sin duplicados y ordenados}.
En C\texttt{++} se implementa mediante \texttt{set}, normalmente como
un árbol balanceado.

A diferencia de \texttt{unordered\_set}, los elementos siempre se
mantienen ordenados según el comparador utilizado.

\textbf{Complejidad:} insertar, buscar y borrar $O(\log n)$.
Consultar el menor o mayor elemento mediante \texttt{begin()} o
\texttt{rbegin()} es $O(1)$.

\begin{lstlisting}[style=cpp]
set<int> s;                    // CREATE: conjunto vacio

s.insert(5);                   // INSERT: O(log n)
s.insert(1);                   // INSERT: O(log n)
s.insert(3);                   // INSERT: O(log n)
s.insert(3);                   // no se agrega: ya existe

int menor = *s.begin();        // READ: minimo -> O(1)
int mayor = *s.rbegin();       // READ: maximo -> O(1)

auto it = s.find(3);           // SEARCH: O(log n)
if (it != s.end())
    cout << *it;

s.erase(3);                    // DELETE: O(log n)

int n = s.size();              // READ: cantidad -> O(1)

// Recorrer en orden ascendente
for (int x : s)
    cout << x << " ";
// O(n)
\end{lstlisting}

\textbf{lower\_bound y upper\_bound:}

\begin{lstlisting}[style=cpp]
set<int> s = {1, 3, 5, 7, 9};

auto it1 = s.lower_bound(5);
// primer elemento >= 5
// resultado: 5

auto it2 = s.upper_bound(5);
// primer elemento > 5
// resultado: 7
\end{lstlisting}

Esto permite resolver consultas de \textbf{piso y techo} rápidamente:

\begin{lstlisting}[style=cpp]
// FLOOR: mayor elemento <= x
auto it = s.upper_bound(x);

if (it != s.begin()) {
    --it;
    cout << *it;
}

// CEILING: menor elemento >= x
auto it2 = s.lower_bound(x);

if (it2 != s.end())
    cout << *it2;
\end{lstlisting}

\textbf{Cuándo usarlo:}
\begin{itemize}
    \item Cuando se necesitan elementos únicos y ordenados.
    \item Obtener mínimo o máximo dinámicamente.
    \item Encontrar el siguiente elemento mayor o igual.
    \item Encontrar el anterior elemento menor o igual.
    \item Mantener un conjunto mientras se insertan y eliminan elementos.
\end{itemize}

\textbf{Diferencia clave:} si solo necesitas saber si existe un elemento,
\texttt{unordered\_set} suele ser más rápido. Si necesitas consultar
orden, \texttt{lower\_bound}, \texttt{upper\_bound}, mínimo o máximo,
\texttt{set} es la opción adecuada.


\subsection{Mapas ordenados (TreeMap)}

Mapa \textbf{clave $\rightarrow$ valor} cuyas claves se mantienen
ordenadas. En C\texttt{++} se implementa mediante \texttt{map}, usando
un árbol balanceado.

Cada clave aparece una sola vez. Los valores pueden repetirse.

\textbf{Complejidad:} insertar, buscar, acceder y borrar por clave
$O(\log n)$. \texttt{begin()} permite obtener la menor clave en
$O(1)$.

\begin{lstlisting}[style=cpp]
map<string,int> m;             // CREATE: mapa vacio

m["b"] = 2;                    // INSERT: O(log n)
m["a"] = 1;                    // INSERT: O(log n)
m["c"] = 3;                    // INSERT: O(log n)

int v = m["a"];                // READ: O(log n)
// v = 1

m["a"] = 10;                   // UPDATE: O(log n)

m.erase("b");                  // DELETE: O(log n)

auto it = m.find("c");         // SEARCH: O(log n)
if (it != m.end())
    cout << it->second;

// Primera pareja: menor clave
auto primera = m.begin();      // O(1)
cout << primera->first;

// Ultima pareja: mayor clave
auto ultima = prev(m.end());   // O(1) amortizado
cout << ultima->first;
\end{lstlisting}

\textbf{Las claves se recorren ordenadas:}

\begin{lstlisting}[style=cpp]
map<int,string> m;

m[30] = "C";
m[10] = "A";
m[20] = "B";

for (auto [clave, valor] : m)
    cout << clave << " " << valor;

// Salida:
// 10 A
// 20 B
// 30 C
\end{lstlisting}

\textbf{lower\_bound y upper\_bound} también funcionan sobre las claves:

\begin{lstlisting}[style=cpp]
map<int,string> m;
m[10] = "A";
m[20] = "B";
m[30] = "C";

auto it = m.lower_bound(15);
// primera clave >= 15
// apunta a 20

auto it2 = m.upper_bound(20);
// primera clave > 20
// apunta a 30
\end{lstlisting}

\textbf{Uso clásico: contar frecuencias ordenadas.}

\begin{lstlisting}[style=cpp]
map<int,int> freq;

for (int x : datos)
    freq[x]++;

// Las claves aparecen en orden creciente.
for (auto [x, cantidad] : freq)
    cout << x << ": " << cantidad << "\n";
\end{lstlisting}

\textbf{Cuándo usarlo:}
\begin{itemize}
    \item Cuando se necesitan pares clave--valor ordenados.
    \item Contar frecuencias y procesarlas en orden.
    \item Encontrar claves inmediatamente mayores o menores.
    \item Procesar eventos o valores manteniendo un orden dinámico.
    \item Cuando se necesitan operaciones de árbol como
    \texttt{lower\_bound} y \texttt{upper\_bound}.
\end{itemize}


\subsection*{Comparación rápida}

\begin{center}
\small
\begin{tabular}{|l|c|c|c|}
\hline
\textbf{Estructura} & \textbf{Duplicados} & \textbf{Orden} & \textbf{Buscar} \\
\hline
\texttt{unordered\_set} & No & No & $O(1)$ prom. \\
\hline
\texttt{unordered\_map} & Claves: No & No & $O(1)$ prom. \\
\hline
\texttt{set} & No & Sí & $O(\log n)$ \\
\hline
\texttt{map} & Claves: No & Sí & $O(\log n)$ \\
\hline
\end{tabular}
\end{center}

\textbf{Regla práctica:}

\begin{itemize}[nosep,leftmargin=*]
    \item \texttt{unordered\_set} $\rightarrow$ \textit{¿Ya existe este elemento?}
    \item \texttt{unordered\_map} $\rightarrow$ \textit{¿Qué valor corresponde a esta clave?}
    \item \texttt{set} $\rightarrow$ \textit{Elementos únicos + ordenados.}
    \item \texttt{map} $\rightarrow$ \textit{Claves $\rightarrow$ valores + claves ordenadas.}
\end{itemize}


\section{Trees}

\section{Range Queries}

\subsection{Casos de uso}

\begin{table}[H]
\raggedright
\scriptsize
\setlength{\tabcolsep}{2pt}
\renewcommand{\arraystretch}{1.1}
\begin{tabular}{|p{0.43\columnwidth}|p{0.29\columnwidth}|p{0.20\columnwidth}|}
\hline
\textbf{Cuando el problema te pida:} &
\textbf{Entonces usa:} &
\textbf{Complejidad} \\
\hline

Suma de rango sobre un arreglo que \textbf{nunca cambia}, con muchas consultas
(ej. ``suma de ventas entre el día $l$ y $r$'') &
Prefix Sums &
Build $O(n)$ \newline Query $O(1)$ \\
\hline

Suma de rango con actualizaciones \textbf{puntuales} frecuentes
(ej. ``suma el stock de la posición $i$, luego consulta un rango'') &
Fenwick Tree (BIT) &
Build $O(n)$ \newline Query/Update $O(\log n)$ \\
\hline

Sumar $v$ a \textbf{todo un rango} $[l,r]$ y también consultar la
\textbf{suma} de cualquier rango &
FenwickRangeUpdate (doble BIT) &
Query/Update $O(\log n)$ \\
\hline

Range update + range query con operaciones más complejas:
min/max en rango, asignar un valor a todo un rango, contar ceros, etc. &
Segment Tree + Lazy Propagation &
Query/Update $O(\log n)$ \\
\hline

Mínimo, máximo o gcd de cualquier rango sobre un arreglo
\textbf{estático} (ej. RMQ clásico, LCA) &
Sparse Table &
Build $O(n\log n)$ \\ Query $O(1)$ \\
\hline

Mínimo o máximo de rango cuando el arreglo \textbf{sí cambia}
con el tiempo &
Segment Tree &
Query/Update $O(\log n)$ \\
\hline

Consultas 2D: cuántos puntos o cuánta suma hay dentro de un
rectángulo $[x_1,x_2]\times[y_1,y_2]$ &
Fenwick 2D &
Query/Update \newline $O(\log n\log m)$ \\
\hline

Encontrar el $k$-ésimo elemento, o ``menor índice cuya suma acumulada
alcanza $X$'' (frecuencias) &
Fenwick con \texttt{lowerBound} &
$O(\log n)$ \\
\hline

Muchas queries de rango \textbf{offline} (las conoces todas de antemano)
sobre un arreglo estático o con pocos updates &
Mo's Algorithm &
$O((n+q)\sqrt{n})$ \\
\hline

LCA entre dos nodos, o mínimo en rango sobre un recorrido Euler de un
árbol (valores consecutivos difieren en $\pm1$) &
RMQ $\pm1$ \newline (Euler Tour + Sparse Table) &
Build $O(n)$ \newline Query $O(1)$ \\
\hline

Cuántos elementos en el rango $[l,r]$ están dentro de un rango de
\textbf{valores} $[a,b]$, o cuál es el $k$-ésimo menor en $[l,r]$ &
Wavelet Tree &
Query $O(\log n)$ \\
\hline

Queries o updates sobre \textbf{caminos o subárboles} de un árbol
(ej. ``suma el camino de $u$ a $v$'') &
Heavy-Light Decomposition \newline (+ Segment/Fenwick) &
Query/Update \newline $O(\log^2 n)$ \\
\hline

Necesitas algo simple de implementar bajo presión, con balance
razonable entre update y query, sin memorizar mucha teoría &
Sqrt Decomposition &
Query/Update $O(\sqrt{n})$ \\
\hline

\end{tabular}
\caption{Guía rápida: qué estructura de range queries usar según el problema, con su complejidad.}
\end{table}



\subsection{Prefix Sums}

\subsection{Suma de Euler Totient ($\phi$) en un rango}

\textbf{Idea:} $\phi(n)$ cuenta cuántos enteros en $[1,n]$ son
coprimos con $n$. Calcular $\phi$ de \emph{un} número es teoría de
números ($O(\sqrt n)$ factorizando); pero cuando el problema pide la
\textbf{suma de $\phi(i)$ para $i$ en un rango $[l,r]$}, con muchas
queries, el patrón es en dos fases: (1) precalculas el arreglo de
$\phi$ una sola vez con una criba, y (2) respondes las queries de
rango con las estructuras de esta sección (Prefix Sums o Fenwick)
sobre ese arreglo ya calculado. La fase (1) es teoría de números; la
fase (2) es exactamente Range Queries.

\textbf{¿Por qué usarlo?} Factorizar cada número del rango uno por
uno para responder cada query sería $O(q\sqrt n)$ y no escala.
Precalcular $\phi$ una vez y sumar con prefijos reduce cada query a
$O(1)$, sin importar cuántas se hagan.

\textbf{Casos de uso:}

\begin{itemize}
    \item Suma de $\phi(i)$ (u otra función aritmética calculable por
    criba, como $\mu$ o el número de divisores) en un rango $[l,r]$,
    con $l,r \le N$ y $N$ manejable ($\le 10^7$ aprox.).
    \item El mismo problema pero con $[L,R]$ muy grandes
    ($R \le 10^{12}$) y $R-L$ pequeño: usa criba segmentada en vez de
    criba desde 1.
    \item Cualquier variante donde $\phi$ no cambia entre queries
    (arreglo estático) $\Rightarrow$ Prefix Sums. Si además necesitas
    modificar valores puntuales entre queries, usa Fenwick.
\end{itemize}

\textbf{Complejidad:} criba lineal $O(N)$, criba segmentada
$O((R-L+1)\log\log R + \sqrt R)$; en ambos casos, una vez armado el
arreglo, sumar un rango es $O(1)$ con Prefix Sums o $O(\log N)$ con
Fenwick si hace falta actualizar.


\subsubsection{Criba lineal de $\phi$ + Prefix Sums (rango $[1,N]$)}

\begin{lstlisting}[style=cpp]
struct PhiRangeSum {
    vector<long long> phi, pref;
    int n;

    // CREATE: calcula phi[1..n] con criba lineal y sus prefijos.
    // phi[i] = cantidad de enteros en [1, i] coprimos con i.
    PhiRangeSum(int n) : n(n) {
        phi.assign(n + 1, 0);
        pref.assign(n + 1, 0);

        vector<int> primos;
        vector<bool> compuesto(n + 1, false);
        phi[1] = 1;

        for (int i = 2; i <= n; i++) {
            if (!compuesto[i]) {
                primos.push_back(i);
                phi[i] = i - 1;
            }

            for (int p : primos) {
                if ((long long)i * p > n) break;

                compuesto[i * p] = true;

                if (i % p == 0) {
                    phi[i * p] = phi[i] * p;
                    break;
                } else {
                    phi[i * p] = phi[i] * (p - 1);
                }
            }
        }

        for (int i = 1; i <= n; i++)
            pref[i] = pref[i - 1] + phi[i];
    }                                          // O(n)

    // READ: suma de phi(i) para i en [l, r].
    long long rango(int l, int r) {
        if (l > r) return 0;
        return pref[r] - pref[l - 1];
    }                                          // O(1)

    // READ: valor individual de phi(i).
    long long get(int i) {
        return phi[i];
    }                                          // O(1)
};
\end{lstlisting}

\textbf{Ejemplo de funcionamiento:}

\begin{lstlisting}[style=cpp]
PhiRangeSum pr(10);
// phi(1..10) = [1, 1, 2, 2, 4, 2, 6, 4, 6, 4]

cout << pr.get(7);
// phi(7) = 6  (7 es primo)

cout << pr.rango(1, 10);
// 1+1+2+2+4+2+6+4+6+4 = 32

cout << pr.rango(4, 8);
// 2+4+2+6+4 = 18
\end{lstlisting}


\subsubsection{Criba segmentada de $\phi$ (rango $[L,R]$ grande)}

\textbf{Idea:} si $R$ es muy grande ($10^{12}$) pero $R-L$ es
manejable, no puedes cribar desde 1. En vez de eso, precalculas los
primos hasta $\sqrt R$ con una criba normal, y para cada número del
segmento $[L,R]$ le vas quitando sus factores primos pequeños,
aplicando la fórmula $\phi(n) = n\prod_{p\mid n}(1-\tfrac1p)$. Si al
final queda un residuo $>1$, es un único primo grande (mayor a
$\sqrt R$) que también se aplica a la fórmula.

\begin{lstlisting}[style=cpp]
struct PhiSegmentedRangeSum {
    vector<long long> phi, pref;
    long long L, R;

    // CREATE: calcula phi(i) para cada i en [L, R] factorizando con
    // los primos hasta sqrt(R), y sus prefijos sobre el segmento.
    PhiSegmentedRangeSum(long long L, long long R) : L(L), R(R) {
        int lim = (int)sqrt((double)R) + 1;

        vector<int> primos;
        vector<bool> compuesto(lim + 1, false);
        for (int i = 2; i <= lim; i++) {
            if (!compuesto[i]) primos.push_back(i);
            for (int p : primos) {
                if ((long long)i * p > lim) break;
                compuesto[i * p] = true;
                if (i % p == 0) break;
            }
        }                                      // O(sqrt(R) log log sqrt(R))

        int sz = R - L + 1;
        vector<long long> num(sz);
        phi.assign(sz, 0);

        for (int i = 0; i < sz; i++) {
            num[i] = L + i;
            phi[i] = L + i;
        }

        for (int p : primos) {
            long long start = ((L + p - 1) / p) * p;

            for (long long j = start; j <= R; j += p) {
                int idx = j - L;
                phi[idx] -= phi[idx] / p;

                while (num[idx] % p == 0)
                    num[idx] /= p;
            }
        }

        // lo que quede > 1 es un primo grande (> sqrt(R)).
        for (int i = 0; i < sz; i++)
            if (num[i] > 1)
                phi[i] -= phi[i] / num[i];

        pref.assign(sz + 1, 0);
        for (int i = 0; i < sz; i++)
            pref[i + 1] = pref[i] + phi[i];
    }                                          // O((R-L+1) log log R + sqrt(R))

    // READ: suma de phi(i) para i en [l, r], con L <= l <= r <= R.
    long long rango(long long l, long long r) {
        return pref[r - L + 1] - pref[l - L];
    }                                          // O(1)

    // READ: valor individual de phi(i), con L <= i <= R.
    long long get(long long i) {
        return phi[i - L];
    }                                          // O(1)
};
\end{lstlisting}

\textbf{Ejemplo de funcionamiento:}

\begin{lstlisting}[style=cpp]
PhiSegmentedRangeSum ps(95, 105);
// calcula phi(95), phi(96), ..., phi(105)
// sin cribar desde 1

cout << ps.get(100);
// phi(100) = 40

cout << ps.rango(95, 105);
// suma de phi(95) a phi(105)

cout << ps.rango(100, 102);
// phi(100) + phi(101) + phi(102)
\end{lstlisting}

\textbf{Cuál variante usar:}

\begin{table}[H]
\centering
\tiny
\setlength{\tabcolsep}{1.5pt}
\renewcommand{\arraystretch}{1.1}
\begin{tabular}{|p{0.32\columnwidth}|p{0.30\columnwidth}|p{0.27\columnwidth}|}
\hline
\textbf{Situación} & \textbf{Estructura} & \textbf{Complejidad} \\
\hline
$[1,N]$, $N$ chico/mediano
& \texttt{PhiRangeSum}
& $O(N)$ \\
\hline
$[L,R]$ grande, $R-L$ chico
& \texttt{PhiSegmentedRangeSum}
& $O((R{-}L)\log\log R+\sqrt{R})$ \\
\hline
$\phi$ cambia entre queries
& Fenwick sobre $\phi$
& $O(\log N)$ \\
\hline
\end{tabular}
\caption{Variantes para la suma de $\phi$.}
\label{tab:phi_variants}
\end{table}


\textbf{Regla práctica:} el cálculo de $\phi$ en sí (la criba, lineal
o segmentada) es teoría de números y no cambia según la query. Una vez
que tienes el arreglo de $\phi$ listo, tratarlo como \emph{cualquier
otro arreglo}: si no cambia, Prefix Sums; si necesitas actualizarlo
puntualmente entre queries, Fenwick. No hace falta una estructura
especial para $\phi$ más allá de construir bien el arreglo inicial.

\subsection{Árbol de Fenwick (BIT)}

\textbf{Idea:} un arreglo normal da suma de prefijo en $O(1)$ pero
actualizar un elemento cuesta $O(n)$ (o al revés si guardas prefijos).
El Fenwick equilibra ambas a $O(\log n)$ guardando \emph{sumas
parciales inteligentes}: la posición $i$ almacena la suma de un bloque
de tamaño $i\,\&\,(-i)$ (su bit menos significativo) que termina en
$i$. Para consultar un prefijo saltas hacia atrás quitando el bit bajo
($i\mathrel{-}=i\,\&\,(-i)$); para actualizar subes sumándolo
($i\mathrel{+}=i\,\&\,(-i)$). Cada camino toca a lo más $\log n$
posiciones (un bit por paso). Es 1-indexado.

\textbf{¿Por qué usarlo?} Es mucho más simple y liviano que un Segment
Tree cuando solo necesitas sumas (o cualquier operación invertible)
de prefijo/rango con actualizaciones puntuales. Menos memoria, menos
código y una constante más pequeña.

\textbf{Casos de uso:}

\begin{itemize}
    \item Suma de prefijo / suma de rango con updates frecuentes.
    \item Conteo de inversiones en un arreglo.
    \item Frecuencias acumuladas y búsqueda del $k$-ésimo elemento.
    \item Range update + range query cuando no hace falta un Segment
    Tree con Lazy Propagation.
    \item Consultas 2D: conteo/suma de puntos dentro de un rectángulo.
\end{itemize}

\textbf{Complejidad:} construir $O(n)$, update puntual $O(\log n)$,
query de prefijo/rango $O(\log n)$, range update + range query
$O(\log n)$ y versión 2D $O(\log n \cdot \log m)$.


\subsubsection{BIT clásico (Point Update, Range Query)}

\begin{lstlisting}[style=cpp]
struct FenwickTree {
    vector<long long> t;
    int n;

    // CREATE: reserva un arbol de tamano n (1-indexado).
    FenwickTree(int n) : t(n + 1, 0), n(n) {}

    // CREATE: construye el arbol a partir de un arreglo a[0..n-1]
    // propagando cada valor a su "padre" una sola vez, en vez de
    // hacer n updates (que serian O(n log n)).
    void build(vector<long long>& a) {
        for (int i = 1; i <= n; i++) {
            t[i] += a[i - 1];

            int j = i + (i & (-i));

            if (j <= n)
                t[j] += t[i];
        }
    }                                      // O(n)

    // UPDATE: suma v al elemento en la posicion i.
    void update(int i, long long v) {
        for (; i <= n; i += i & (-i))
            t[i] += v;
    }                                      // O(log n)

    // UPDATE: fija el valor absoluto de la posicion i en x.
    void set(int i, long long x) {
        update(i, x - get(i));
    }                                      // O(log n)

    // READ: suma de prefijo [1..i].
    long long query(int i) {
        long long s = 0;

        for (; i > 0; i -= i & (-i))
            s += t[i];

        return s;
    }                                      // O(log n)

    // READ: suma del rango [l, r].
    long long rango(int l, int r) {
        if (l > r)
            return 0;

        return query(r) - query(l - 1);
    }                                      // O(log n)

    // READ: valor puntual en la posicion i.
    long long get(int i) {
        return rango(i, i);
    }                                      // O(log n)

    // SEARCH: menor indice i tal que query(i) >= objetivo.
    // Usa binary lifting sobre el arbol.
    // Requiere que todos los valores sean >= 0.
    int lowerBound(long long objetivo) {
        int pos = 0;
        long long ac = 0;

        int LOG = 0;

        while ((1 << LOG) <= n)
            LOG++;

        for (int k = LOG - 1; k >= 0; k--) {
            int nx = pos + (1 << k);

            if (nx <= n && ac + t[nx] < objetivo) {
                pos = nx;
                ac += t[nx];
            }
        }

        return pos + 1;
    }                                      // O(log n)

    // DELETE: reinicia el arbol a ceros.
    void clear() {
        fill(t.begin(), t.end(), 0);
    }                                      // O(n)
};
\end{lstlisting}


\textbf{Ejemplo de funcionamiento:}

\begin{lstlisting}[style=cpp]
vector<long long> a = {5, 2, 9, 1, 7};

FenwickTree bit(5);
bit.build(a);

// arreglo logico:
// [5, 2, 9, 1, 7]
// posiciones 1..5

cout << bit.query(3);
// 5 + 2 + 9 = 16

cout << bit.rango(2, 4);
// 2 + 9 + 1 = 12

bit.update(2, 3);

// arreglo logico:
// [5, 5, 9, 1, 7]

cout << bit.get(2);
// 5

bit.set(4, 10);

// arreglo logico:
// [5, 5, 9, 10, 7]

cout << bit.lowerBound(15);
// menor indice cuyo prefijo alcanza 15
// -> indice 3 (5 + 5 + 9 = 19 >= 15)
\end{lstlisting}


\subsubsection{BIT con Range Update + Range Query (Doble BIT)}

\textbf{Idea:} con \emph{dos} Fenwick y el truco de diferencias se
puede sumar $v$ a todo un rango $[l,r]$ en $O(\log n)$ y también
consultar la suma de cualquier rango en $O(\log n)$. Esto permite
resolver simultáneamente actualizaciones y consultas de rango.

\begin{lstlisting}[style=cpp]
struct FenwickRangeUpdate {
    FenwickTree b1, b2;
    int n;

    // CREATE: reserva dos arboles auxiliares de tamano n.
    FenwickRangeUpdate(int n) : b1(n), b2(n), n(n) {}

    // UPDATE: suma v a todos los elementos del rango [l, r].
    void rangoUpdate(int l, int r, long long v) {
        b1.update(l, v);
        b1.update(r + 1, -v);

        b2.update(l, v * (l - 1));
        b2.update(r + 1, -v * r);
    }                                      // O(log n)

    // READ: suma de prefijo [1..i].
    long long prefijo(int i) {
        return b1.query(i) * i - b2.query(i);
    }                                      // O(log n)

    // READ: suma del rango [l, r].
    long long rango(int l, int r) {
        if (l > r)
            return 0;

        return prefijo(r) - prefijo(l - 1);
    }                                      // O(log n)

    // READ: valor puntual en la posicion i.
    long long get(int i) {
        return rango(i, i);
    }                                      // O(log n)
};
\end{lstlisting}

\textbf{Ejemplo de funcionamiento:}

\begin{lstlisting}[style=cpp]
FenwickRangeUpdate f(5);
// arreglo logico inicial: [0, 0, 0, 0, 0]

f.rangoUpdate(2, 4, 3);
// arreglo logico: [0, 3, 3, 3, 0]

cout << f.rango(1, 5);
// 0+3+3+3+0 = 9

f.rangoUpdate(1, 3, 2);
// arreglo logico: [2, 5, 5, 3, 0]

cout << f.get(2);
// 5
\end{lstlisting}

\subsubsection{BIT 2D (point update, sum de submatriz)}

\textbf{Idea:} un BIT dentro de otro BIT. Cada update/query en $x$
dispara un update/query completo en $y$, dando $O(\log n \cdot
\log m)$ en vez de $O(nm)$ por consulta.

\begin{lstlisting}[style=cpp]
struct Fenwick2D {
    vector<vector<long long>> t;
    int n, m;

    // CREATE: reserva una matriz de arboles de tamano n x m.
    Fenwick2D(int n, int m) : t(n + 1, vector<long long>(m + 1, 0)),
                               n(n), m(m) {}

    // UPDATE: suma v a la celda (x, y).
    void update(int x, int y, long long v) {
        for (int i = x; i <= n; i += i & (-i))
            for (int j = y; j <= m; j += j & (-j))
                t[i][j] += v;
    }                                          // O(log n * log m)

    // READ: suma de la submatriz [1..x] x [1..y].
    long long query(int x, int y) {
        long long s = 0;
        for (int i = x; i > 0; i -= i & (-i))
            for (int j = y; j > 0; j -= j & (-j))
                s += t[i][j];
        return s;
    }                                          // O(log n * log m)

    // READ: suma del rectangulo [x1..x2] x [y1..y2].
    long long rango(int x1, int y1, int x2, int y2) {
        return query(x2, y2) - query(x1 - 1, y2)
             - query(x2, y1 - 1) + query(x1 - 1, y1 - 1);
    }                                          // O(log n * log m)
};
\end{lstlisting}

\textbf{Ejemplo de funcionamiento:}

\begin{lstlisting}[style=cpp]
Fenwick2D f(4, 4);

f.update(1, 1, 5);
f.update(2, 3, 7);
f.update(4, 4, 2);

cout << f.rango(1, 1, 2, 3);
// 5 + 7 = 12

cout << f.rango(3, 3, 4, 4);
// 2
\end{lstlisting}

\textbf{Diferencia con un arreglo de prefijos simple:}

\begin{table}[H]
\raggedright
\scriptsize
\setlength{\tabcolsep}{3pt}
\begin{tabular}{|l|c|c|c|}
\hline
\textbf{Estructura} & \textbf{Construir} & \textbf{Query} & \textbf{Update} \\
\hline
Prefijos simples & $O(n)$ & $O(1)$ & $O(n)$ \\
\hline
\texttt{FenwickTree} & $O(n)$ & $O(\log n)$ & $O(\log n)$ \\
\hline
\texttt{FenwickRangeUpdate} & $O(n)$ & $O(\log n)$ & $O(\log n)$* \\
\hline
\texttt{Fenwick2D} & $O(nm)$ & $O(\log n\log m)$ & $O(\log n\log m)$ \\
\hline
\end{tabular}
\caption{Comparación de complejidades: Fenwick vs. prefijos simples.}
\end{table}

\noindent
*Range update + range query, con la variante de doble BIT.

\textbf{Regla práctica:} si el arreglo cambia con frecuencia y solo
necesitas sumas (o cualquier operación invertible) de prefijo o rango,
el Fenwick da el mejor balance entre simplicidad de código y
velocidad. Si necesitas min/max, operaciones no invertibles, o updates
de rango con lazy propagation general, usa Segment Tree en su lugar.



\subsection{Sqrt Decomposition}

\subsection{Sparse Table (RMQ)}

\textbf{Idea:} responde consultas de rango en $O(1)$ sobre un arreglo
que \textbf{no cambia}, precalculando \texttt{sp[k][i]} $=$ el
resultado de combinar el bloque que empieza en $i$ y tiene largo
$2^k$. La clave está en la consulta: para el rango $[l,r]$ tomas
$k=\lfloor\log_2(r-l+1)\rfloor$ y lo cubres con \emph{dos} bloques de
largo $2^k$ —uno pegado a $l$ y otro pegado a $r$— que
\textbf{se solapan}. Si la operación es \emph{idempotente}
($f(x,x)=x$, como min, max o gcd), contar el solape dos veces no
cambia el resultado, así que la consulta son solo dos lecturas y una
combinación. Para operaciones \textbf{no idempotentes} (como suma o
xor) hace falta la variante \emph{Disjoint Sparse Table}, que evita el
solape por construcción.

\textbf{¿Por qué usarlo?} Es la estructura más rápida posible para
consultas de rango ($O(1)$ por query) cuando el arreglo es estático.
A cambio de esa velocidad, no soporta actualizaciones: reconstruir
cuesta $O(n\log n)$, así que si el arreglo cambia, conviene un
Segment Tree.

\textbf{Casos de uso:}

\begin{itemize}
    \item RMQ clásico: mínimo o máximo de cualquier rango, arreglo
    estático.
    \item gcd o AND/OR de rango sobre un arreglo estático.
    \item LCA vía Euler Tour + RMQ $\pm1$.
    \item Suma, xor o cualquier operación no idempotente de rango
    estático, en $O(1)$, usando Disjoint Sparse Table.
    \item Cualquier problema con muchísimas queries offline sobre datos
    que nunca se modifican.
\end{itemize}

\textbf{Complejidad:} construir $O(n\log n)$, consulta $O(1)$. (Si el
arreglo \emph{cambia}, usa Segment Tree: $O(\log n)$ por operación.)


\subsubsection{Sparse Table clásico (operaciones idempotentes: min, max, gcd)}

\begin{lstlisting}[style=cpp]
struct SparseTable {
    vector<vector<long long>> sp;
    vector<int> lg;
    int n;

    // CREATE: precalcula la tabla a partir de un arreglo a[0..n-1].
    // sp[k][i] = minimo del bloque [i, i + 2^k - 1].
    SparseTable(vector<long long>& a) {
        n = a.size();

        lg.assign(n + 1, 0);
        for (int i = 2; i <= n; i++)
            lg[i] = lg[i / 2] + 1;

        int K = lg[n] + 1;
        sp.assign(K, vector<long long>(n));
        sp[0] = a;

        for (int k = 1; k < K; k++)
            for (int i = 0; i + (1 << k) <= n; i++)
                sp[k][i] = min(sp[k - 1][i],
                                sp[k - 1][i + (1 << (k - 1))]);
    }                                          // O(n log n)

    // READ: minimo del rango [l, r] (0-indexado, inclusivo).
    // Cubre [l, r] con dos bloques de largo 2^k que se solapan;
    // como min es idempotente, el solape no afecta el resultado.
    long long query(int l, int r) {
        int k = lg[r - l + 1];
        return min(sp[k][l], sp[k][r - (1 << k) + 1]);
    }                                          // O(1)
};
\end{lstlisting}

\textbf{Ejemplo de funcionamiento:}

\begin{lstlisting}[style=cpp]
vector<long long> a = {5, 2, 9, 1, 7, 3};

SparseTable st(a);
// arreglo logico: [5, 2, 9, 1, 7, 3]  (posiciones 0..5)

cout << st.query(0, 2);
// min(5, 2, 9) = 2

cout << st.query(2, 5);
// min(9, 1, 7, 3) = 1

cout << st.query(4, 4);
// min(7) = 7
\end{lstlisting}

\textbf{Nota:} para \texttt{max} o \texttt{gcd} solo cambia
\texttt{min(...)} por \texttt{max(...)} o \texttt{\_\_gcd(...)} en el
constructor y en \texttt{query}. Si necesitas saber \emph{qué índice}
alcanza el mínimo (no solo el valor), guarda índices en vez de
valores y compara \texttt{a[sp[k][i]]} en lugar de \texttt{sp[k][i]}.


\subsubsection{Disjoint Sparse Table (operaciones no idempotentes: suma, xor)}

\textbf{Idea:} en vez de dos bloques que se solapan, cada nivel $k$
particiona el arreglo en bloques de largo $2^{k+1}$ y precalcula, para
cada bloque, un prefijo hacia la izquierda desde su punto medio y un
prefijo hacia la derecha desde ese mismo punto medio. Toda consulta
$[l,r]$ cae exactamente en un nivel donde $l$ y $r$ quedan en mitades
distintas del mismo bloque (el nivel lo da el bit más alto donde
difieren $l$ y $r$), así que se combina sin solapar y sirve para
operaciones \textbf{no idempotentes} como la suma.

\begin{lstlisting}[style=cpp]
struct DisjointSparseTable {
    vector<vector<long long>> pre;
    vector<long long> a;
    int n, LOG;

    // CREATE: precalcula prefijos izquierda/derecha por nivel,
    // partiendo el arreglo en bloques de tamano 2^(k+1).
    DisjointSparseTable(vector<long long>& arr) {
        a = arr;
        n = a.size();
        LOG = 1;
        while ((1 << LOG) < n) LOG++;
        LOG++;

        pre.assign(LOG, vector<long long>(n));

        for (int k = 0; k < LOG; k++) {
            int block = 1 << (k + 1);

            for (int start = 0; start < n; start += block) {
                int mid = min(start + block / 2, n);
                int end = min(start + block, n);

                if (mid > n) continue;

                // prefijo hacia la izquierda desde mid-1
                if (mid - 1 >= start) {
                    pre[k][mid - 1] = a[mid - 1];
                    for (int i = mid - 2; i >= start; i--)
                        pre[k][i] = a[i] + pre[k][i + 1];
                }

                // prefijo hacia la derecha desde mid
                if (mid < end) {
                    pre[k][mid] = a[mid];
                    for (int i = mid + 1; i < end; i++)
                        pre[k][i] = pre[k][i - 1] + a[i];
                }
            }
        }
    }                                          // O(n log n)

    // READ: suma del rango [l, r] (0-indexado, inclusivo).
    // El nivel k lo da el bit mas alto donde l y r difieren:
    // ahi es donde el bloque se parte en dos mitades disjuntas.
    long long query(int l, int r) {
        if (l == r) return a[l];

        int k = 31 - __builtin_clz(l ^ r);
        return pre[k][l] + pre[k][r];
    }                                          // O(1)
};
\end{lstlisting}

\textbf{Ejemplo de funcionamiento:}

\begin{lstlisting}[style=cpp]
vector<long long> a = {5, 2, 9, 1, 7, 3};

DisjointSparseTable dst(a);
// arreglo logico: [5, 2, 9, 1, 7, 3]  (posiciones 0..5)

cout << dst.query(0, 2);
// 5 + 2 + 9 = 16

cout << dst.query(2, 5);
// 9 + 1 + 7 + 3 = 20

cout << dst.query(1, 1);
// 2
\end{lstlisting}

\textbf{Diferencia entre las dos variantes:}

\begin{table}[H]
\raggedright
\scriptsize
\setlength{\tabcolsep}{3pt}
\begin{tabular}{|l|c|c|c|}
\hline
\textbf{Estructura} & \textbf{Construir} & \textbf{Query} & \textbf{Updates} \\
\hline
Sparse Table (min/max/gcd) & $O(n\log n)$ & $O(1)$ & No \\
\hline
Disjoint Sparse Table (suma/xor) & $O(n\log n)$ & $O(1)$ & No \\
\hline
Segment Tree & $O(n)$ & $O(\log n)$ & Sí, $O(\log n)$ \\
\hline
\end{tabular}
\caption{Sparse Table vs. Disjoint Sparse Table vs. Segment Tree.}
\end{table}

\textbf{Regla práctica:} si el arreglo \textbf{no cambia}, usa Sparse
Table para min/max/gcd, o Disjoint Sparse Table si necesitas suma,
xor, o cualquier operación no idempotente. En cuanto el arreglo
\textbf{cambia} —aunque sea una sola actualización— reconstruir en
$O(n\log n)$ deja de valer la pena y conviene un Segment Tree.



\subsection{Segment Tree}

\subsection{RMQ (Range Minimum Query)}

\subsection{Mo's Algorithm}

\subsection{Wavelet Tree}

\subsection{Heavy-Light Decomposition}


%---------------------------------------------------------------------------
\part{Grafos y Árboles}
%---------------------------------------------------------------------------

\section{Grafos (Graphs)}
\pendiente

\section{Árboles (Trees)}
\pendiente

\section{Disjoint Set Union (DSU)}
\pendiente


%---------------------------------------------------------------------------
\part{Ordenamiento y Búsqueda}
%---------------------------------------------------------------------------

\section{Ordenamiento y Búsqueda (Sorting \& Searching)}
Varias búsquedas lineales y sub-lineales. (La búsqueda binaria tiene
su propia sección.)

\subsection{Número faltante (Missing Number)}
\textbf{Idea:} los números $0,1,\dots,n$ deberían sumar
$\frac{n(n+1)}{2}$ (fórmula de Gauss). Si falta uno, la suma real queda
corta \emph{exactamente} en ese valor, así que el faltante es
(suma esperada $-$ suma real). Basta un recorrido para acumular la suma.
Alternativa sin riesgo de overflow: hacer XOR de $0..n$ y de todos los
elementos; lo que quede es el faltante. \textbf{Complejidad:} $O(n)$.
\begin{lstlisting}[style=cpp]
long long esp = (long long)n*(n+1)/2;          // suma 0..n
long long falta = esp - accumulate(a.begin(),a.end(),0LL);
\end{lstlisting}


\subsection{Máximo / Mínimo (Array Max/Min)}
\textbf{Idea:} recorres el arreglo una sola vez llevando el menor y el
mayor vistos hasta ahora; comparás cada elemento contra ambos y
actualizas. No requiere ordenar (ordenar sería $O(n\log n)$, peor). Con
un recorrido consigues los dos a la vez. \textbf{Complejidad:} $O(n)$.
\begin{lstlisting}[style=cpp]
int mn = *min_element(a.begin(), a.end());
int mx = *max_element(a.begin(), a.end());
\end{lstlisting}


\subsection{Par de suma absoluta mínima}
\textbf{Idea:} buscamos el par cuya suma esté \emph{más cerca de cero}.
Tras ordenar, ponemos un puntero en cada extremo. La suma de los
extremos te dice hacia dónde moverte: si es negativa, para acercarla a
$0$ hay que subir el menor (\texttt{lo++}); si es positiva, hay que
bajar el mayor (\texttt{hi--}). Cada movimiento descarta el candidato
menos prometedor, y en el camino guardas el $|suma|$ más pequeño visto.
Los dos punteros recorren el arreglo una vez ($O(n)$); domina el sort.
\textbf{Complejidad:} $O(n\log n)$.
\begin{lstlisting}[style=cpp]
sort(a.begin(), a.end());
int lo=0, hi=a.size()-1; long long mej=LLONG_MAX;
while(lo<hi){ long long s=(long long)a[lo]+a[hi];
  mej=min(mej,llabs(s));
  if(s<0) lo++; else hi--; }
\end{lstlisting}


\subsection{Par con diferencia dada (Difference Pair)}
\textbf{Idea:} ¿existe un par cuya diferencia sea exactamente $d$? Tras
ordenar, dos punteros avanzan hacia adelante: si la diferencia actual
$a_j-a_i$ es menor que $d$, adelantas el puntero lejano ($j$) para
agrandarla; si es mayor, adelantas el cercano ($i$) para achicarla; si
es igual, encontraste el par. Como el arreglo está ordenado, la
diferencia crece con $j$ y decrece con $i$, así que el barrido nunca
retrocede. \textbf{Complejidad:} $O(n\log n)$ (domina el sort).
\begin{lstlisting}[style=cpp]
sort(a.begin(),a.end()); int i=0,j=1,n=a.size();
while(i<n && j<n){
  if(i!=j && a[j]-a[i]==d) return true;
  if(a[j]-a[i]<d) j++; else i++; }
\end{lstlisting}


\subsection{Búsqueda ternaria (Ternary Search)}
\textbf{Idea:} sirve para hallar el pico de una función
\textbf{unimodal} (sube y luego baja, con un solo máximo — no para buscar
un valor en un arreglo cualquiera). Partes el rango en tres con dos
puntos $m_1<m_2$ y comparas: si $f(m_1)<f(m_2)$, el máximo no puede estar
a la izquierda de $m_1$, así que descartas ese tercio; en caso contrario
descartas el tercio derecho. Cada paso elimina $\sim\frac{1}{3}$ del
rango, por eso converge en logarítmico. (Para un mínimo, invierte la
comparación.) \textbf{Complejidad:} $O(\log n)$ — NO $O(n\log n)$.
\begin{lstlisting}[style=cpp]
while(hi-lo>2){ long long m1=lo+(hi-lo)/3, m2=hi-(hi-lo)/3;
  if(f(m1)<f(m2)) lo=m1+1; else hi=m2-1; }
// revisa [lo,hi] y devuelve el x que maximiza f
\end{lstlisting}


\subsection{Búsqueda de Fibonacci (Fibonacci Search)}
\textbf{Idea:} misma estrategia que la binaria (descartar mitades de un
arreglo ordenado), pero en vez de partir por el punto medio elige el
punto de comparación usando números de Fibonacci consecutivos. Mantiene
tres Fibonacci y en cada paso mira la posición \texttt{off+f2}; según el
elemento sea menor o mayor que el objetivo, retrocede uno o dos pasos en
la secuencia. Ventaja: solo usa \emph{sumas y restas}, nunca división,
lo que la hace útil en hardware sin división rápida o para saltar por
posiciones amigables con la caché. \textbf{Complejidad:} $O(\log n)$.
\begin{lstlisting}[style=cpp]
int f2=0,f1=1,fb=f1+f2;
while(fb<n){ f2=f1; f1=fb; fb=f1+f2; }
int off=-1;
while(fb>1){ int i=min(off+f2,n-1);
  if(a[i]<x){ fb=f1; f1=f2; f2=fb-f1; off=i; }
  else if(a[i]>x){ fb=f2; f1=f1-f2; f2=fb-f1; }
  else return i; }
if(f1==1 && off+1<n && a[off+1]==x) return off+1;
\end{lstlisting}

(En Java es idéntico al C\texttt{++}, cambiando \texttt{a.length}.)

\subsection{Búsqueda exponencial (Exponential Search)}
\textbf{Idea:} dos fases sobre un arreglo \textbf{ordenado}. Primero
\emph{acotas} el objetivo doblando el índice ($1,2,4,8,\dots$) hasta que
$a[i]$ lo supera; en ese momento sabes que el objetivo está entre $i/2$
e $i$. Segundo, haces una binaria normal dentro de ese bloque. Como el
bloque encontrado tiene tamaño proporcional a la posición del objetivo,
es ideal cuando el arreglo es enorme o de tamaño desconocido y el
objetivo está cerca del inicio. \textbf{Complejidad:} $O(\log n)$.
\begin{lstlisting}[style=cpp]
if(a[0]==x) return 0;
int i=1; while(i<n && a[i]<=x) i*=2;   // dobla el rango
// luego binaria en [i/2, min(i,n-1)]
\end{lstlisting}

\subsection{Búsqueda por saltos (Jump Search)}
\textbf{Idea:} en vez de revisar elemento por elemento, avanzas en
bloques de tamaño $\sqrt{n}$ mirando solo el \emph{último} de cada
bloque; cuando ese último ya supera al objetivo, sabes que (si existe)
está en el bloque anterior, así que lo barres linealmente. El balance
$\sqrt{n}$ saltos $+$ $\sqrt{n}$ del barrido minimiza el total. Punto
medio entre la lineal ($O(n)$) y la binaria ($O(\log n)$), útil cuando
saltar hacia atrás es costoso. \textbf{Complejidad:} $O(\sqrt{n})$.
\begin{lstlisting}[style=cpp]
int paso=max(1,(int)sqrt((double)n)), prev=0;
while(prev<n && a[min(prev+paso,n)-1]<x) prev+=paso;
for(int i=prev; i<min(prev+paso,n); i++)
  if(a[i]==x) return i;
\end{lstlisting}


\subsection{Heap Sort}
\textbf{Idea:} un \emph{heap} binario se representa en el mismo arreglo:
el hijo izquierdo de $i$ es $2i+1$ y el derecho $2i+2$; en un
\emph{max-heap} todo padre es $\geq$ sus hijos, así que el mayor está en
la raíz. Heap Sort tiene dos fases. \textbf{(1) Construir el heap}:
recorres de abajo hacia arriba aplicando \texttt{bajar} (sift-down) a
cada padre. Aunque parezca $O(n\log n)$, es $O(n)$: casi todos los nodos
están cerca de las hojas y se hunden poco (la suma
$\sum n/2^h\cdot h$ converge a $2n$). \textbf{(2) Ordenar}: repetidamente
intercambias la raíz (el mayor) con el último elemento activo —quedando
ya ordenado al final— y ``hundes'' la nueva raíz para restaurar el heap;
cada una de las $n$ extracciones cuesta $O(\log n)$. In-place, sin
memoria extra, pero \textbf{no estable}. \textbf{Complejidad:}
$O(n\log n)$ en todos los casos.
\begin{lstlisting}[style=cpp]
void bajar(vector<int>&a,int n,int i){    // hunde: O(log n)
  while(true){ int m=i,l=2*i+1,r=2*i+2;
    if(l<n&&a[l]>a[m]) m=l;
    if(r<n&&a[r]>a[m]) m=r;
    if(m==i) break; swap(a[i],a[m]); i=m; } }
void heapSort(vector<int>&a){ int n=a.size();
  for(int i=n/2-1;i>=0;i--) bajar(a,n,i);   // heap: O(n)
  for(int i=n-1;i>0;i--){ swap(a[0],a[i]);  // max al final
    bajar(a,i,0); } }
\end{lstlisting}

(En Java es idéntico al C\texttt{++} con \texttt{a.length}.)

\subsection{QuickSelect}
\textbf{Idea:} encuentra el $k$-ésimo menor \emph{sin ordenar todo}.
Como quicksort, particiona alrededor de un pivote dejando a la izquierda
lo menor y a la derecha lo mayor; pero como sabes en qué lado cae la
posición $k$, recurres (o iteras) sobre \emph{un solo} lado y descartas
el otro. Por eso el promedio baja de $O(n\log n)$ a $O(n)$: cada fase
procesa \emph{la mitad} del anterior, y $n+\tfrac{n}{2}+\tfrac{n}{4}+
\dots = 2n$. El peor caso $O(n^2)$ aparece con pivotes pésimos (arreglo
casi ordenado + pivote extremo); \textbf{elegir el pivote al azar} lo
vuelve prácticamente imposible. La versión de abajo es iterativa (sin
pila de recursión). Si necesitas garantía $O(n)$ en el peor caso, existe
\emph{mediana de medianas} para elegir un buen pivote.
\textbf{Complejidad:} promedio $O(n)$, peor $O(n^2)$.
\begin{lstlisting}[style=cpp]
int quickSelect(vector<int> a, int k){       // k in [0,n-1]
  int lo=0, hi=a.size()-1;
  while(lo<hi){
    int p=a[lo+rand()%(hi-lo+1)], i=lo, j=hi; // pivote aleatorio
    while(i<=j){ while(a[i]<p) i++; while(a[j]>p) j--;
      if(i<=j){ swap(a[i],a[j]); i++; j--; } }
    if(k<=j) hi=j;                             // k a la izq
    else if(k>=i) lo=i;                        // k a la der
    else return a[k];
  }
  return a[lo];
}
// C++ trae nth_element(a.begin(), a.begin()+k, a.end()); // O(n)
\end{lstlisting}

(En Java igual, con \texttt{Random}; código completo en el repo.)

\section{Búsqueda Binaria (Binary Search)}
\textbf{Idea central:} mantienes un rango $[lo,hi]$ donde \emph{sabes}
que está la respuesta y en cada paso lo partes por la mitad, descartando
la mitad que no puede contenerla. Como reduces a la mitad cada vez,
llegas en $\log_2 n$ pasos. Funciona si los datos están \textbf{ordenados}
o, más general, si existe un \textbf{predicado monótono} (una condición
que pasa de falsa a verdadera una sola vez): entonces buscas la frontera
del cambio. \textbf{Complejidad:} $O(\log n)$ en discreto,
$O(\text{iter})$ en reales, $O(L\log n)$ con cadenas. Tip: usa
\texttt{lo+(hi-lo)/2} para el \texttt{mid} y evitar overflow.

\subsection{En arreglo ordenado}
\textbf{Cómo:} comparas $x$ con el elemento del medio; si es menor,
descartas la mitad derecha, si es mayor la izquierda. \texttt{lower\_bound}
te da el primer índice con valor $\geq x$ y \texttt{upper\_bound} el
primer $> x$; con eso sabes si $x$ existe y cuántas copias hay.
\begin{lstlisting}[style=cpp]
int i = lower_bound(a.begin(),a.end(),x)-a.begin(); // >= x
int j = upper_bound(a.begin(),a.end(),x)-a.begin(); // >  x
bool hay = binary_search(a.begin(),a.end(),x);
\end{lstlisting}


\subsection{Menor $x$ que cumple}
\textbf{Cómo:} aquí no buscas en un arreglo sino en el \emph{rango de
respuestas posibles}. El predicado es monótono F..F\,T..T (una vez que
algo ``sirve'', todo lo mayor también). En cada paso pruebas el
\texttt{mid}: si cumple, lo guardas como candidato y sigues buscando a la
izquierda (\texttt{hi=mid-1}) por si hay uno menor que también sirva; si
no cumple, subes (\texttt{lo=mid+1}). Al final \texttt{res} es la
frontera izquierda: el \emph{menor} que cumple. Ejemplo: ``mínima
capacidad para repartir en $\leq k$ días''.
\begin{lstlisting}[style=cpp]
long long res=-1;
while(lo<=hi){ long long mid=lo+(hi-lo)/2;
  if(cumple(mid)){res=mid; hi=mid-1;} // sirve -> izq
  else lo=mid+1; }
\end{lstlisting}


\subsection{Mayor $x$ que cumple}
\textbf{Cómo:} el caso espejo. Ahora el predicado es T..T\,F..F (sirve
hasta cierto punto y luego deja de servir). Cuando el \texttt{mid}
cumple, lo guardas y buscas a la \emph{derecha} (\texttt{lo=mid+1}) por
si hay uno mayor que también sirva; si no cumple, bajas
(\texttt{hi=mid-1}). Respecto al anterior solo cambian esas dos líneas.
Ejemplo: ``máxima longitud de corte tal que salgan $\geq k$ piezas''.
\begin{lstlisting}[style=cpp]
long long res=-1;
while(lo<=hi){ long long mid=lo+(hi-lo)/2;
  if(cumple(mid)){res=mid; lo=mid+1;} // sirve -> der
  else hi=mid-1; }
\end{lstlisting}


\subsection{Sobre reales}
\textbf{Cómo:} cuando la respuesta es un número \emph{real} (raíces,
puntos de equilibrio) no hay ``mitad entera'' ni condición $lo\leq hi$
para parar. En su lugar iteras un número fijo de veces (unas 100 basta:
el rango se divide por $2^{100}$, precisión más que suficiente),
moviendo \texttt{lo} o \texttt{hi} al \texttt{mid} según el predicado.
Al terminar, \texttt{lo}$\approx$\texttt{hi} es la frontera buscada.
\begin{lstlisting}[style=cpp]
for(int it=0; it<100; it++){ double mid=(lo+hi)/2;
  if(cumple(mid)) hi=mid; else lo=mid; }
// lo ~ hi = frontera
\end{lstlisting}


\subsection{Con strings}
\textbf{Cómo:} exactamente la misma binaria, pero comparando cadenas en
orden \emph{lexicográfico} (diccionario) en lugar de números. La única
diferencia de costo es que comparar dos cadenas cuesta $O(L)$ (longitud),
no $O(1)$, así que el total sube a $O(L\log n)$. Ojo: el arreglo debe
estar ordenado con el mismo criterio (ASCII: mayúsculas antes que
minúsculas).
\begin{lstlisting}[style=cpp]
// vector<string> s ordenado
int i = lower_bound(s.begin(),s.end(),string("caro"))
        - s.begin();
bool hay = binary_search(s.begin(),s.end(),string("caro"));
\end{lstlisting}


%---------------------------------------------------------------------------
\part{Matemáticas}
%---------------------------------------------------------------------------

\section{Matemáticas (Math)}
\pendiente

\section{Teoría de Números (Number Theory)}
\pendiente

\section{Combinatoria (Combinatorics)}
\pendiente

\section{Máscaras de Bits (Bitmasking)}
\pendiente

\section{Geometría (Geometry)}
\pendiente

%---------------------------------------------------------------------------
\part{Técnicas Algorítmicas}
%---------------------------------------------------------------------------

\section{Programación Dinámica (Dynamic Programming)}
\pendiente

\section{Algoritmos Voraces (Greedy)}
\pendiente

\section{Ad-hoc y Simulación}
\pendiente

%---------------------------------------------------------------------------
\part{Cadenas}
%---------------------------------------------------------------------------

\section{Cadenas (Strings)}
\pendiente

\end{document}
